%%%%%%%%%%%%%%%%%%%%%%%%%%%%%%%%%%%%%%%%%%%%%%%%%%%%%%%%%%%%%%%%%%%%%%%%%%%%%%%%
%  IEEE Conference Paper — Card Fraud Triage
%  Project 26 — CSAI AML & DL, Newton School of Technology / Rishihood University
%%%%%%%%%%%%%%%%%%%%%%%%%%%%%%%%%%%%%%%%%%%%%%%%%%%%%%%%%%%%%%%%%%%%%%%%%%%%%%%%
\documentclass[conference,10pt]{IEEEtran}
\IEEEoverridecommandlockouts

%%% ---- packages ----
\usepackage{cite}
\usepackage{amsmath,amssymb,amsfonts}
\usepackage{algorithmic}
\usepackage{graphicx}
\usepackage{textcomp}
\usepackage{booktabs}
\usepackage{multirow}
\usepackage{url}
\usepackage{array}
\usepackage{balance}
\usepackage{xcolor}
\usepackage{enumitem}
\usepackage{subcaption}
\usepackage[hyphens]{url}

\graphicspath{{figures/}}

\def\BibTeX{{\rm B\kern-.05em{\sc i\kern-.025em b}\kern-.08em
    T\kern-.1667em\lower.7ex\hbox{E}\kern-.125emX}}

\begin{document}

%%%%%%%%%%%%%%%%%%%%%%%%%%%%%%%%%%%%%%%%%%%%%%%%%%%%%%%%%%%%%%%%%%%%%%%%%%%%%%%%
%  TITLE & AUTHORS
%%%%%%%%%%%%%%%%%%%%%%%%%%%%%%%%%%%%%%%%%%%%%%%%%%%%%%%%%%%%%%%%%%%%%%%%%%%%%%%%
\title{Card Fraud Triage via Gaussian Mixture Density Scoring,\\
Deep Autoencoders, and Hybrid Latent-Space Modeling}

\author{
\IEEEauthorblockN{%
Ayush Sahu\IEEEauthorrefmark{1},
Rachit Singh\IEEEauthorrefmark{2},
Ashar Asif Shaikh\IEEEauthorrefmark{3}}
\IEEEauthorblockA{%
Department of Computer Science and Artificial Intelligence\\
Newton School of Technology (NST) / Rishihood University, Sonipat, India\\
\texttt{\{ayush.sahu, rachit.singh, ashar.shaikh\}@nst.rishihood.edu.in}}
}

\maketitle

%%%%%%%%%%%%%%%%%%%%%%%%%%%%%%%%%%%%%%%%%%%%%%%%%%%%%%%%%%%%%%%%%%%%%%%%%%%%%%%%
%  ABSTRACT
%%%%%%%%%%%%%%%%%%%%%%%%%%%%%%%%%%%%%%%%%%%%%%%%%%%%%%%%%%%%%%%%%%%%%%%%%%%%%%%%
\begin{abstract}
Credit-card fraud detection presents a challenging machine-learning problem
characterised by extreme class imbalance ($<\!0.2\%$ fraud prevalence),
high-dimensional feature spaces, and non-stationary adversarial patterns.
Conventional supervised classifiers trained on historical fraud labels are
susceptible to concept drift, label noise, and catastrophic overfitting to the
majority class.  This paper reframes fraud detection as an \emph{unsupervised
anomaly-detection and triage} problem.  We train models \emph{exclusively on
legitimate transactions} and evaluate them through three progressively
sophisticated paradigms:
%
(1)~a \textbf{38-component diagonal-covariance Gaussian Mixture Model (GMM)}
that estimates the probability density of normal spending and uses negative
log-likelihood as the anomaly score (Average Precision, AP~=~0.2719);
%
(2)~an \textbf{Undercomplete Denoising Autoencoder} with a 6-dimensional
latent bottleneck that learns non-linear manifold representations and scores
anomalies via per-sample reconstruction error (AP~=~0.7248, $+166\%$);
%
(3)~a \textbf{Hybrid Latent-GMM Detector} that fits an 8-component GMM
directly in the autoencoder's compressed latent space and fuses reconstruction
error with latent-density scores, achieving the best overall performance
(AP~=~\textbf{0.7612}, $+180\%$ over baseline).
%
All experiments are conducted on the publicly available European credit-card
transaction dataset (283,726 cleaned transactions).  A chronological
70/15/15 train--validation--test split is enforced to prevent temporal data
leakage, and feature standardisation is fitted solely on legitimate training
records.  Systematic ablation studies on bottleneck dimensionality, denoising
noise level, and activation function validate every architectural choice.
Operational triage evaluation under realistic banking constraints---including
false-positive-rate budgets ($\le$\,0.1\%, 0.5\%, 1.0\%) and fixed daily
human-review capacities (Top-50, 100, 200)---demonstrates that the hybrid
model identifies \textbf{38 of 52 holdout frauds with only 9 false alarms},
achieving 76.0\% precision in a Top-50 review queue.  These findings confirm
that deep non-linear bottleneck representations combined with latent density
fusion substantially overcome the geometric limitations of convex Gaussian
clustering in high-dimensional financial transaction triage.
\end{abstract}

\begin{IEEEkeywords}
credit-card fraud detection, anomaly detection, unsupervised learning,
Gaussian mixture model, autoencoder, hybrid modeling, deep learning,
class imbalance, transaction triage, ablation study
\end{IEEEkeywords}

%%%%%%%%%%%%%%%%%%%%%%%%%%%%%%%%%%%%%%%%%%%%%%%%%%%%%%%%%%%%%%%%%%%%%%%%%%%%%%%%
%  I. INTRODUCTION
%%%%%%%%%%%%%%%%%%%%%%%%%%%%%%%%%%%%%%%%%%%%%%%%%%%%%%%%%%%%%%%%%%%%%%%%%%%%%%%%
\section{Introduction}
\label{sec:intro}

Electronic payment fraud inflicts global losses exceeding \$32 billion
annually~\cite{nilson2023}, requiring real-time automated detection systems
that balance sensitivity against the operational cost of false alarms.
Financial institutions typically employ \emph{triage pipelines}: automated
scoring ranks every transaction by suspicion, and a fixed-capacity team of
human analysts investigates only the top-$N$ flagged cases each day.

Supervised fraud classifiers---Random Forests, Gradient-Boosted Trees,
feed-forward neural networks---have demonstrated strong in-distribution
performance~\cite{he2009}.  However, they suffer from three structural
limitations in production settings:

\begin{enumerate}[leftmargin=*]
\item \textbf{Severe class imbalance.}  Fraud constitutes $<\!0.2\%$ of
      transactions, skewing decision boundaries and inflating false-negative
      rates unless aggressive resampling or cost-sensitive losses are applied.
\item \textbf{Concept drift.}  Adversarial actors continuously evolve fraud
      tactics, rendering historical fraud labels stale within weeks.
\item \textbf{Label noise.}  Ground-truth annotations depend on delayed
      chargebacks and manual investigations, introducing systematic labelling
      errors.
\end{enumerate}

Unsupervised anomaly detection sidesteps these issues by learning a model of
\emph{normality} from the abundant legitimate-transaction population.
Transactions that deviate significantly from the learned distribution are
flagged for investigation without requiring any fraud labels during training.

This paper presents a systematic, multi-paradigm investigation:

\begin{itemize}[leftmargin=*]
\item \textbf{Phase~1 (AML Baseline):} A classical Gaussian Mixture Model
      (GMM) performs parametric density estimation in the original 30-dimensional
      feature space, establishing the machine-learning baseline.
\item \textbf{Phase~2 (Deep Learning):} An Undercomplete Denoising Autoencoder
      compresses transactions through a non-linear bottleneck; per-sample
      reconstruction error serves as the anomaly score.
\item \textbf{Phase~2+ (Hybrid Innovation):} A novel Hybrid Latent-GMM
      Detector fits a compact GMM directly in the autoencoder's latent space,
      then fuses reconstruction error with latent-density scores through a
      weighted linear combination.
\end{itemize}

All models are evaluated under a leakage-free experimental protocol with
chronological data splitting and standardisation fitted exclusively on
legitimate training records.

The remainder of this paper is organised as follows.
Section~\ref{sec:related} reviews related work.
Section~\ref{sec:data} describes the dataset and preprocessing pipeline.
Section~\ref{sec:method} details the three modeling paradigms.
Section~\ref{sec:results} reports experimental results, ablation studies, and
operational triage benchmarks.
Section~\ref{sec:discussion} analyses findings and limitations.
Section~\ref{sec:conclusion} concludes the paper.

%%%%%%%%%%%%%%%%%%%%%%%%%%%%%%%%%%%%%%%%%%%%%%%%%%%%%%%%%%%%%%%%%%%%%%%%%%%%%%%%
%  II. RELATED WORK
%%%%%%%%%%%%%%%%%%%%%%%%%%%%%%%%%%%%%%%%%%%%%%%%%%%%%%%%%%%%%%%%%%%%%%%%%%%%%%%%
\section{Related Work}
\label{sec:related}

\subsection{Anomaly Detection}
Chandola et al.~\cite{chandola2009} survey anomaly-detection methods across
statistical, distance-based, and machine-learning categories.  Probabilistic
density models such as Gaussian Mixture Models~\cite{dempster1977} and kernel
density estimators have long been applied to network-intrusion and fraud
detection.  However, parametric models struggle in high-dimensional spaces
where covariance matrices become ill-conditioned or singular.

\subsection{Autoencoders for Anomaly Detection}
Sakurada and Yairi~\cite{sakurada2014} demonstrated that undercomplete
autoencoders can detect anomalies via reconstruction error when trained solely
on normal data.  Vincent et al.~\cite{vincent2010} introduced Denoising
Autoencoders (DAE) that inject controlled noise during training, forcing the
encoder to learn more robust latent representations.  Bank et al.~\cite{bank2023}
provide a modern survey of autoencoder architectures and training strategies
for representation learning.

\subsection{Hybrid and Deep GMM Models}
Zong et al.~\cite{zong2018} proposed the Deep Autoencoding Gaussian Mixture
Model (DAGMM), which jointly trains a compression network and a GMM estimation
network in an end-to-end fashion.  Our approach differs by using a
\emph{two-stage} pipeline---pre-training the autoencoder, then fitting a
classical GMM in the frozen latent space---offering greater interpretability
and simpler hyperparameter tuning at the cost of joint optimality.

\subsection{Fraud Detection under Imbalance}
He and Garcia~\cite{he2009} review algorithmic strategies for imbalanced
learning, including cost-sensitive learning, oversampling (SMOTE), and
ensemble methods.  Dal~Pozzolo et al.~\cite{ulb_dataset} introduced the
European credit-card dataset used in this study and benchmarked calibration
methods under severe class imbalance.

%%%%%%%%%%%%%%%%%%%%%%%%%%%%%%%%%%%%%%%%%%%%%%%%%%%%%%%%%%%%%%%%%%%%%%%%%%%%%%%%
%  III. DATASET AND PREPROCESSING
%%%%%%%%%%%%%%%%%%%%%%%%%%%%%%%%%%%%%%%%%%%%%%%%%%%%%%%%%%%%%%%%%%%%%%%%%%%%%%%%
\section{Dataset and Preprocessing}
\label{sec:data}

\subsection{Dataset Description}
We use the publicly available European credit-card fraud detection
dataset~\cite{ulb_dataset}, which contains 284,807 anonymised transactions
recorded over a 48-hour period in September~2013.  Each transaction is
represented by 30 features: elapsed time since the first transaction
(\texttt{Time}), 28 principal-component features (\texttt{V1}--\texttt{V28})
obtained via a confidential PCA transformation, and the monetary
\texttt{Amount}.  The binary label \texttt{Class} takes the value~0 for
legitimate transactions and~1 for fraud.  The dataset contains 284,315
legitimate and 492 fraudulent transactions (0.1727\% base rate).

\subsection{Data Cleaning}
Exploratory data analysis revealed 1,081 exact-duplicate rows.  These were
removed to prevent artificial metric inflation and information leakage during
cross-validation, yielding a cleaned dataset of 283,726 transactions with 473
remaining fraud cases (0.1667\% base rate).

\subsection{Chronological Partitioning}
To simulate realistic deployment conditions and prevent temporal leakage,
transactions were sorted by \texttt{Time} and split into three contiguous
partitions as summarised in Table~\ref{tab:split}.

\begin{table}[htbp]
\caption{Chronological Data Partitioning After Cleaning}
\label{tab:split}
\centering
\small
\begin{tabular}{@{}lrrrr@{}}
\toprule
Partition & Total & Legitimate & Fraud & Fraud Rate \\
\midrule
Training (70\%)   & 198,608 & 198,242 & 366 & 0.184\% \\
Validation (15\%) &  42,559 &  42,504 &  55 & 0.129\% \\
Test (15\%)       &  42,559 &  42,507 &  52 & 0.122\% \\
\midrule
\textbf{Total}    & 283,726 & 283,253 & 473 & 0.167\% \\
\bottomrule
\end{tabular}
\end{table}

\subsection{Feature Standardisation}
A \texttt{StandardScaler} was fitted \emph{exclusively} on the 198,242
legitimate training transactions.  Fraud cases in the training partition were
withheld from both scaler fitting and model training to maintain the
unsupervised setup.  The fitted scaler parameters ($\hat{\mu}_j$,
$\hat{\sigma}_j$ for each feature $j$) were then applied to transform
validation and test features:
\begin{equation}
\tilde{x}_j = \frac{x_j - \hat{\mu}_j}{\hat{\sigma}_j}
\end{equation}

This procedure ensures that no information from fraud transactions or future
time steps contaminates the training process.

%%%%%%%%%%%%%%%%%%%%%%%%%%%%%%%%%%%%%%%%%%%%%%%%%%%%%%%%%%%%%%%%%%%%%%%%%%%%%%%%
%  IV. METHODOLOGY
%%%%%%%%%%%%%%%%%%%%%%%%%%%%%%%%%%%%%%%%%%%%%%%%%%%%%%%%%%%%%%%%%%%%%%%%%%%%%%%%
\section{Methodology}
\label{sec:method}

\subsection{Phase 1: Gaussian Mixture Density Scoring}
\label{subsec:gmm}

A Gaussian Mixture Model (GMM) represents the probability density of
legitimate transaction vectors $\mathbf{x} \in \mathbb{R}^{30}$ as a weighted
sum of $K$ multivariate Gaussian components:
\begin{equation}
p(\mathbf{x}) = \sum_{k=1}^{K} \pi_k \,
    \mathcal{N}\!\left(\mathbf{x} \mid \boldsymbol{\mu}_k,
    \boldsymbol{\Sigma}_k\right)
\label{eq:gmm}
\end{equation}
where $\pi_k$ are non-negative mixture weights summing to unity,
$\boldsymbol{\mu}_k \in \mathbb{R}^{30}$ are component means, and
$\boldsymbol{\Sigma}_k$ are covariance matrices.  Parameters are estimated
via Expectation--Maximisation (EM)~\cite{dempster1977}.

\subsubsection{Model Selection}
An extensive grid search over $K \in \{1, 2, \ldots, 40\}$ was performed using
three complementary criteria: Akaike Information Criterion (AIC), Bayesian
Information Criterion (BIC), and legitimate-validation log-likelihood.  A
diagonal covariance structure ($\boldsymbol{\Sigma}_k = \mathrm{diag}
(\sigma_{k,1}^2, \ldots, \sigma_{k,30}^2)$) was adopted to avoid
ill-conditioning in 30 dimensions.  Covariance regularisation was applied
($\text{reg\_covar} = 10^{-5}$), and each configuration was initialised three
times ($n_{\text{init}} = 3$) to mitigate sensitivity to random initialisation.

$K = 36$ achieved the highest single-run validation log-likelihood.  However,
$K = 38$ was selected as the final configuration based on superior AIC/BIC
trade-off and greater stability across repeated seeds (see
Fig.~\ref{fig:aic_bic}).  The model converged in 111 EM iterations.

\begin{figure}[t]
\centering
\includegraphics[width=\columnwidth]{aic_bic.png}
\caption{AIC and BIC as a function of GMM component count $K$.  The elbow
region around $K = 36$--$38$ balances model complexity against fit quality.}
\label{fig:aic_bic}
\end{figure}

\subsubsection{Anomaly Scoring}
The anomaly score for a transaction vector $\mathbf{x}$ is defined as the
negative log-likelihood under the fitted GMM:
\begin{equation}
s_{\text{GMM}}(\mathbf{x}) = -\log p(\mathbf{x})
\label{eq:nll}
\end{equation}
Higher scores indicate that the transaction is less probable under the learned
legitimate distribution and thus more suspicious.

\subsubsection{Threshold Selection}
The binary decision threshold $\tau^*$ was chosen to maximise the F1-score on
the validation set:
\begin{equation}
\tau^* = \arg\max_{\tau} \; F_1(\tau) = \arg\max_{\tau} \;
    \frac{2 \cdot P(\tau) \cdot R(\tau)}{P(\tau) + R(\tau)}
\end{equation}
The optimal validation threshold was $\tau^* = 94.7575$.

\subsection{Phase 2: Undercomplete Denoising Autoencoder}
\label{subsec:ae}

\subsubsection{Architecture}
The autoencoder consists of a symmetric encoder--decoder pair:

\smallskip
\noindent\textbf{Encoder:}
$\mathbb{R}^{30} \xrightarrow{\text{Linear}} \mathbb{R}^{20}
\xrightarrow{\text{BN+LeakyReLU+Drop}} \mathbb{R}^{20}
\xrightarrow{\text{Linear}} \mathbb{R}^{12}
\xrightarrow{\text{BN+LeakyReLU+Drop}} \mathbb{R}^{12}
\xrightarrow{\text{Linear}} \mathbb{R}^{6}$

\smallskip
\noindent\textbf{Decoder:}
$\mathbb{R}^{6} \xrightarrow{\text{Linear}} \mathbb{R}^{12}
\xrightarrow{\text{BN+LeakyReLU+Drop}} \mathbb{R}^{12}
\xrightarrow{\text{Linear}} \mathbb{R}^{20}
\xrightarrow{\text{BN+LeakyReLU+Drop}} \mathbb{R}^{20}
\xrightarrow{\text{Linear}} \mathbb{R}^{30}$

\smallskip
\noindent
The network contains 1,826 trainable parameters.  Batch Normalization (BN)
stabilises intermediate activations, LeakyReLU ($\alpha = 0.1$) prevents dying
neurons on zero-mean standardised inputs, and Dropout ($p = 0.05$) provides
lightweight regularisation.

\subsubsection{Denoising Training}
During each training forward pass, additive Gaussian noise is injected into the
input:
\begin{equation}
\tilde{\mathbf{x}} = \mathbf{x} + \boldsymbol{\epsilon}, \quad
\boldsymbol{\epsilon} \sim \mathcal{N}(\mathbf{0}, \sigma^2 \mathbf{I})
\label{eq:noise}
\end{equation}
with noise factor $\sigma = 0.02$.  The model learns to reconstruct the
\emph{clean} input $\mathbf{x}$ from the corrupted observation
$\tilde{\mathbf{x}}$, forcing the encoder to capture the underlying data
manifold rather than memorising individual samples~\cite{vincent2010}.

\subsubsection{Training Configuration}
The network was trained for 50 epochs using the AdamW optimiser with learning
rate $\eta = 10^{-3}$, weight decay $\lambda = 10^{-5}$, and a
\texttt{ReduceLROnPlateau} scheduler (factor 0.5, patience 5).  Early stopping
with patience~10 on validation loss selected epoch~28 as the best checkpoint
(best validation MSE~=~0.1843).  Training was performed exclusively on the
198,242 legitimate training transactions.

\subsubsection{Anomaly Scoring}
The per-sample reconstruction error serves as the anomaly score:
\begin{equation}
s_{\text{AE}}(\mathbf{x}) = \frac{1}{D} \sum_{j=1}^{D}
    \left(x_j - \hat{x}_j\right)^2
\label{eq:recon}
\end{equation}
where $D = 30$ is the feature dimensionality and $\hat{\mathbf{x}} =
g_\phi(f_\theta(\mathbf{x}))$ is the reconstructed output.  Legitimate
transactions that conform to the learned manifold produce low reconstruction
error, while anomalous transactions deviate substantially.

\subsection{Phase 2+: Hybrid Latent-GMM Detector}
\label{subsec:hybrid}

\subsubsection{Motivation}
The Phase~1 GMM suffers from the \emph{curse of dimensionality} in
$\mathbb{R}^{30}$: diagonal covariance matrices cannot capture inter-feature
correlations, and high-dimensional density estimation is statistically
inefficient.  Conversely, the Phase~2 autoencoder lacks a principled
probabilistic interpretation---reconstruction error alone may not distinguish
between transactions that are genuinely anomalous and those that simply lie
in under-represented but legitimate regions of the data manifold.

The Hybrid Latent-GMM Detector addresses both limitations by combining the
complementary strengths of the two paradigms.

\subsubsection{Architecture}
The hybrid detection pipeline operates in three stages:
\begin{enumerate}[leftmargin=*]
\item \textbf{Latent projection.}  The pre-trained autoencoder encoder maps
      each 30D transaction into the 6D latent manifold:
      $\mathbf{z} = f_\theta(\mathbf{x}) \in \mathbb{R}^6$.
\item \textbf{Latent density estimation.}  An 8-component full-covariance GMM
      is fitted on the legitimate-training latent embeddings
      $\{\mathbf{z}_i\}_{i=1}^{N_{\text{legit}}}$.  In 6 dimensions,
      full covariance estimation is well-conditioned and captures correlations
      that diagonal models miss.
\item \textbf{Score fusion.}  The hybrid anomaly score is a weighted linear
      combination of the normalised (min--max scaled) reconstruction error and
      latent NLL:
      \begin{equation}
      s_{\text{hybrid}}(\mathbf{x}) = \alpha \cdot
          \tilde{s}_{\text{AE}}(\mathbf{x}) + (1 - \alpha) \cdot
          \tilde{s}_{\text{Latent-GMM}}(\mathbf{z})
      \label{eq:hybrid}
      \end{equation}
      where $\tilde{s}$ denotes min--max normalisation to $[0, 1]$ and
      $\alpha = 0.6$ was selected via validation AP maximisation.
\end{enumerate}

\subsubsection{Complementary Anomaly Signals}
\begin{itemize}[leftmargin=*]
\item \textbf{Reconstruction component ($\alpha$-weighted):}  Detects
      \emph{structural violations} where the decoder fails to reconstruct
      unfamiliar spending patterns.
\item \textbf{Latent-density component ($(1\!-\!\alpha)$-weighted):}  Detects
      transactions that project into sparse, low-probability regions of the
      normal latent manifold, even when reconstruction error is moderate.
\end{itemize}

%%%%%%%%%%%%%%%%%%%%%%%%%%%%%%%%%%%%%%%%%%%%%%%%%%%%%%%%%%%%%%%%%%%%%%%%%%%%%%%%
%  V. EXPERIMENTAL RESULTS
%%%%%%%%%%%%%%%%%%%%%%%%%%%%%%%%%%%%%%%%%%%%%%%%%%%%%%%%%%%%%%%%%%%%%%%%%%%%%%%%
\section{Experimental Results}
\label{sec:results}

All results are reported on the held-out chronological \emph{test} partition
(42,559 transactions; 52 frauds; 42,507 legitimate).  Decision thresholds
are determined on the \emph{validation} partition and applied unchanged to
the test set.

\subsection{Tri-Model Classification Benchmark}
Table~\ref{tab:comparison} presents the head-to-head comparison of all three
models at their respective validation-F1-optimal thresholds.

\begin{table}[htbp]
\caption{Test-Set Classification Performance\\(52 Fraud, 42,507 Legitimate Transactions)}
\label{tab:comparison}
\centering
\small
\begin{tabular}{@{}lrrr@{}}
\toprule
Metric & \makecell{Phase~1\\GMM\\($K\!=\!38$)}
       & \makecell{Phase~2\\AE\\($d\!=\!6$)}
       & \makecell{Phase~2+\\Hybrid\\($d\!=\!6,K\!=\!8$)} \\
\midrule
Avg.\ Precision (AP) & 0.2719 & 0.7248 & \textbf{0.7612} \\
Precision             & 46.94\% & 78.72\% & \textbf{80.85\%} \\
Recall                & 44.23\% & 71.15\% & \textbf{73.08\%} \\
F1-Score              & 0.4554 & 0.7475 & \textbf{0.7677} \\
FPR                   & 0.0612\% & 0.0235\% & \textbf{0.0212\%} \\
\midrule
True Positives        & 23 & 37 & \textbf{38} \\
False Positives       & 26 & 10 & \textbf{9} \\
False Negatives       & 29 & 15 & \textbf{14} \\
True Negatives        & 42,481 & 42,497 & \textbf{42,498} \\
\bottomrule
\end{tabular}
\end{table}

The Hybrid model achieves the best performance across \emph{every} metric,
delivering an AP of 0.7612---a \textbf{+180\%} relative improvement over the
GMM baseline and \textbf{+5.0\%} over the standalone autoencoder.  It detects
38 of 52 holdout frauds while generating only 9 false alarms out of 42,507
legitimate transactions.

\begin{figure}[t]
\centering
\includegraphics[width=\columnwidth]{comparative_pr.png}
\caption{Precision--Recall curves for all three models on the chronological
test set.  The dashed line indicates the random no-skill baseline
(fraud prevalence $\approx 0.12\%$).}
\label{fig:pr_curve}
\end{figure}

\subsection{False-Positive-Rate Budget Analysis}
In production systems, the false-positive rate (FPR) is often constrained by
investigation capacity.  Table~\ref{tab:fpr} reports fraud recall at three
operationally relevant FPR budgets.

\begin{table}[htbp]
\caption{Fraud Recall at Controlled FPR Budgets (Test Set)}
\label{tab:fpr}
\centering
\small
\begin{tabular}{@{}lrrr@{}}
\toprule
Target FPR & \makecell{GMM\\Recall / Alarms}
           & \makecell{AE\\Recall / Alarms}
           & \makecell{Hybrid\\Recall / Alarms} \\
\midrule
$\le 1.0\%$ & 73.1\% / 316 & 86.5\% / 345 & \textbf{88.5\%} / 332 \\
$\le 0.5\%$ & 63.5\% / 167 & 80.8\% / 175 & \textbf{82.7\%} / 165 \\
$\le 0.1\%$ & 40.4\% / 19  & 75.0\% / 32  & \textbf{76.9\%} / 29  \\
\bottomrule
\end{tabular}
\end{table}

Under the strictest budget (FPR $\le 0.1\%$), the Hybrid model recovers
\textbf{76.9\%} of all frauds (40 of 52) with only 29 false alarms, compared
to 40.4\% for the GMM and 75.0\% for the standalone AE.

\subsection{Fixed Review-Capacity Analysis}
Table~\ref{tab:review} simulates a triage environment where analysts can
investigate a fixed number of flagged transactions per day.

\begin{table}[htbp]
\caption{Fraud Detection under Fixed Daily Review Capacity}
\label{tab:review}
\centering
\small
\begin{tabular}{@{}lrrr@{}}
\toprule
Top-$N$ & \makecell{GMM\\Found / Prec.}
        & \makecell{AE\\Found / Prec.}
        & \makecell{Hybrid\\Found / Prec.} \\
\midrule
Top 50  & 24 / 48.0\% & 36 / 72.0\% & \textbf{38 / 76.0\%} \\
Top 100 & 28 / 28.0\% & 41 / 41.0\% & \textbf{43 / 43.0\%} \\
Top 200 & 33 / 16.5\% & 46 / 23.0\% & \textbf{47 / 23.5\%} \\
\bottomrule
\end{tabular}
\end{table}

In a \textbf{Top-50 daily review queue}, the Hybrid model identifies 38
confirmed frauds (76.0\% precision) versus 36 for the AE and only 24 for the
GMM.  This directly translates to greater fraud-loss recovery within the same
staffing budget.

\subsection{Ablation Studies}
Systematic ablation experiments were conducted to validate the autoencoder's
key design choices.  Results are summarised in Table~\ref{tab:ablations}.

\begin{table}[htbp]
\caption{Systematic Architectural Ablation Studies (Test Set)}
\label{tab:ablations}
\centering
\small
\begin{tabular}{@{}llrrrr@{}}
\toprule
Factor & Config. & Val MSE & Test AP & Prec. & Rec. \\
\midrule
\multirow{4}{*}{\makecell[l]{Latent\\dim.\ $d$}}
& $d\!=\!3$              & 0.284 & 0.651 & 71.4\% & 67.3\% \\
& $d\!=\!6$ {\scriptsize(ours)} & 0.184 & \textbf{0.725} & \textbf{78.7\%} & \textbf{71.2\%} \\
& $d\!=\!12$             & 0.141 & 0.689 & 73.9\% & 65.4\% \\
& $d\!=\!16$             & 0.119 & 0.612 & 66.7\% & 61.5\% \\
\midrule
\multirow{4}{*}{\makecell[l]{Noise\\$\sigma$}}
& 0.00 {\scriptsize(none)}      & 0.171 & 0.692 & 75.6\% & 65.4\% \\
& 0.02 {\scriptsize(ours)}      & 0.184 & \textbf{0.725} & \textbf{78.7\%} & \textbf{71.2\%} \\
& 0.05                   & 0.219 & 0.708 & 76.1\% & 67.3\% \\
& 0.10                   & 0.310 & 0.634 & 68.9\% & 59.6\% \\
\midrule
\multirow{3}{*}{\makecell[l]{Activ.\\fn.}}
& ReLU                   & --- & 0.684 & 74.5\% & 67.3\% \\
& LeakyReLU {\scriptsize(ours)} & --- & \textbf{0.725} & \textbf{78.7\%} & \textbf{71.2\%} \\
& GELU                   & --- & 0.711 & 77.1\% & 71.2\% \\
\bottomrule
\end{tabular}
\end{table}

\subsubsection{Latent Dimensionality}
Reducing the bottleneck to $d = 3$ causes over-compression (information loss),
degrading test AP to 0.651.  Expanding to $d = 12$ or $d = 16$ allows the
network to begin passing anomalous patterns through the bottleneck
(\emph{identity shortcut}), reducing anomaly separability.  The sweet spot
at $d = 6$ balances compression against information retention.

\subsubsection{Denoising Noise Level}
Without noise injection ($\sigma = 0$), the network overfits to training
samples, reducing generalisation (AP = 0.692).  Mild noise ($\sigma = 0.02$)
acts as an implicit regulariser, forcing the encoder to capture the data
manifold's structure.  Excessive noise ($\sigma = 0.10$) overwhelms the
signal, degrading all metrics.

\subsubsection{Activation Function}
Standard ReLU produces dead neurons on zero-mean standardised inputs, reducing
AP to 0.684.  LeakyReLU ($\alpha = 0.1$) preserves negative gradient flow and
achieves the best results.  GELU provides smooth gating and competitive
performance (AP = 0.711) but at higher computational cost.

%%%%%%%%%%%%%%%%%%%%%%%%%%%%%%%%%%%%%%%%%%%%%%%%%%%%%%%%%%%%%%%%%%%%%%%%%%%%%%%%
%  VI. DISCUSSION
%%%%%%%%%%%%%%%%%%%%%%%%%%%%%%%%%%%%%%%%%%%%%%%%%%%%%%%%%%%%%%%%%%%%%%%%%%%%%%%%
\section{Discussion}
\label{sec:discussion}

\subsection{Why the GMM Underperforms}
The 38-component diagonal GMM achieves an AP of only 0.2719 despite careful
model selection.  In 30 dimensions, diagonal covariance matrices cannot model
inter-feature correlations.  Full covariance estimation is infeasible:
$\frac{30 \times 31}{2} = 465$ free parameters per component $\times$ 38
components would require estimating 17,670 covariance entries, leading to
singular or near-singular matrices.  Consequently, the GMM's density landscape
is too coarse to separate the subtle anomaly signal from normal variation.

\subsection{Autoencoder Advantage}
The autoencoder's non-linear encoder--decoder chain learns a compact 6D
manifold that captures the essential structure of legitimate transactions.
Fraudulent transactions---which exploit unusual combinations of features such
as atypical amounts, merchant categories, and temporal patterns---cannot be
faithfully reconstructed from this bottleneck, producing high reconstruction
error.  The +166\% AP improvement over the GMM demonstrates the value of
non-linear representation learning.

\subsection{Value of Hybrid Fusion}
The Hybrid model's modest but consistent improvement over the standalone AE
(AP: 0.7248 $\rightarrow$ 0.7612, +5.0\%) arises because latent-density
scoring and reconstruction-error scoring capture \emph{complementary} anomaly
signals.  Some fraudulent transactions may reconstruct moderately well but
project into low-density latent regions; conversely, some may reside in dense
latent regions but exhibit high reconstruction error.  Fusing both signals via
Eq.~\eqref{eq:hybrid} provides more robust anomaly ranking.

\subsection{Limitations and Future Work}
\begin{itemize}[leftmargin=*]
\item \textbf{Single dataset:}  Results are reported on one benchmark;
      evaluation on proprietary banking data would strengthen external validity.
\item \textbf{Static model:}  The current pipeline does not adapt to concept
      drift.  Periodic retraining or online learning strategies could be
      explored.
\item \textbf{Interpretability:}  While per-feature reconstruction error
      provides basic explainability, more sophisticated attribution methods
      (SHAP, attention-based decoders) could enhance investigator trust.
\item \textbf{End-to-end training:}  A DAGMM-style~\cite{zong2018} joint
      optimisation of the autoencoder and latent GMM may further improve the
      hybrid model's performance.
\end{itemize}

%%%%%%%%%%%%%%%%%%%%%%%%%%%%%%%%%%%%%%%%%%%%%%%%%%%%%%%%%%%%%%%%%%%%%%%%%%%%%%%%
%  VII. CONCLUSION
%%%%%%%%%%%%%%%%%%%%%%%%%%%%%%%%%%%%%%%%%%%%%%%%%%%%%%%%%%%%%%%%%%%%%%%%%%%%%%%%
\section{Conclusion}
\label{sec:conclusion}

This paper presented a systematic investigation of unsupervised credit-card
fraud triage across three paradigms---classical parametric density estimation,
deep non-linear manifold learning, and hybrid latent-density fusion---under a
leakage-free chronological experimental protocol.

The key findings are:
\begin{enumerate}[leftmargin=*]
\item A 38-component diagonal-covariance GMM establishes a classical baseline
      with AP = 0.2719, limited by high-dimensional covariance instability.
\item An Undercomplete Denoising Autoencoder ($d = 6$) dramatically improves
      fraud detection to AP = 0.7248 (+166\%), demonstrating the power of
      non-linear bottleneck representations.
\item The proposed Hybrid Latent-GMM Detector further improves performance to
      AP = \textbf{0.7612} (+180\% over GMM), detecting \textbf{38 of 52}
      holdout frauds with only \textbf{9 false alarms} and achieving
      \textbf{76.0\%} precision in a Top-50 daily review queue.
\item Systematic ablation studies validate architectural choices:
      $d = 6$ latent dimensionality, $\sigma = 0.02$ denoising noise, and
      LeakyReLU activation each contribute meaningfully to the final
      performance.
\end{enumerate}

These results demonstrate that deep representation learning combined with
latent-space density fusion provides a robust, operationally practical
framework for credit-card fraud triage in severely imbalanced, high-dimensional
transaction environments.

%%%%%%%%%%%%%%%%%%%%%%%%%%%%%%%%%%%%%%%%%%%%%%%%%%%%%%%%%%%%%%%%%%%%%%%%%%%%%%%%
%  REFERENCES
%%%%%%%%%%%%%%%%%%%%%%%%%%%%%%%%%%%%%%%%%%%%%%%%%%%%%%%%%%%%%%%%%%%%%%%%%%%%%%%%
\balance
\begin{thebibliography}{99}

\bibitem{nilson2023}
The Nilson Report, ``Global card fraud losses,'' Issue~1234, 2023.

\bibitem{chandola2009}
V.~Chandola, A.~Banerjee, and V.~Kumar,
``Anomaly detection: A survey,''
\emph{ACM Comput.\ Surv.}, vol.~41, no.~3, pp.~1--58, 2009.

\bibitem{dempster1977}
A.~P.~Dempster, N.~M.~Laird, and D.~B.~Rubin,
``Maximum likelihood from incomplete data via the EM algorithm,''
\emph{J.\ R.\ Stat.\ Soc.\ Ser.\ B}, vol.~39, no.~1, pp.~1--38, 1977.

\bibitem{sakurada2014}
M.~Sakurada and T.~Yairi,
``Anomaly detection using autoencoders with nonlinear dimensionality
reduction,''
in \emph{Proc.\ MLSDA Workshop}, 2014, pp.~4--11.

\bibitem{vincent2010}
P.~Vincent, H.~Larochelle, I.~Lajoie, Y.~Bengio, and P.-A.~Manzagol,
``Stacked denoising autoencoders: Learning useful representations in a deep
network with a local denoising criterion,''
\emph{J.\ Mach.\ Learn.\ Res.}, vol.~11, pp.~3371--3408, 2010.

\bibitem{bank2023}
D.~Bank, N.~Koenigstein, and R.~Giryes,
``Autoencoders,''
in \emph{Machine Learning for Data Science Handbook}, Springer, 2023,
pp.~353--374.

\bibitem{zong2018}
B.~Zong, Q.~Song, M.~R.~Min, W.~Cheng, C.~Lumezanu, D.~Cho, and H.~Chen,
``Deep autoencoding Gaussian mixture model for unsupervised anomaly
detection,''
in \emph{Proc.\ ICLR}, 2018.

\bibitem{he2009}
H.~He and E.~A.~Garcia,
``Learning from imbalanced data,''
\emph{IEEE Trans.\ Knowl.\ Data Eng.}, vol.~21, no.~9, pp.~1263--1284, 2009.

\bibitem{ulb_dataset}
A.~Dal~Pozzolo, O.~Caelen, R.~A.~Johnson, and G.~Bontempi,
``Calibrating probability with undersampling for unbalanced classification,''
in \emph{Proc.\ IEEE SSCI}, 2015, pp.~159--166.

\end{thebibliography}

\end{document}
